\chapter{Lessons}

\section{What the project taught}

The chapters above described decisions in place. This one collects the ones that
were only obvious afterwards, which are the ones worth passing on.

\subsection{Correctness in money is a type decision}

Choosing \code{Decimal} and \code{NUMERIC} at the start removed a whole class of
bugs for the life of the project. No test had to check that a total was off by a
cent, because it cannot be. The equivalent decision made late would have meant
touching every model, every parser and every report.

\subsection{The interface is not the application}

Making the Telegram bot call the same services as the web pages meant that
adding a command was a few lines and never a second implementation. The rule
that route handlers contain no arithmetic is what made that possible, and it is
cheap to follow from day one and expensive to retrofit.

\subsection{Boot time is a feature}

Nobody thinks about import time until a platform suspends the process. Then
several seconds of loading pandas and matplotlib become the reason a bot looks
broken. Moving heavy imports inside functions took an afternoon; a test now
keeps them there.

\subsection{Choose where to fail loudly}

The project swallows errors in one place and propagates them in another, and
both are deliberate. A price that cannot be fetched leaves the old value alone,
because the alternative is a portfolio that stops updating over one delisted
ticker. A backup that cannot be delivered returns \code{500}, because a backup
that fails quietly is worse than none.

\subsection{Free tiers are a design constraint}

Sleeping applications and pausing databases are not bugs to work around once;
they change what the software has to contain. The health endpoints, the split
between an in-process scheduler and HTTP task endpoints, and the backup that
lives outside the provider all exist because of that constraint.

\subsection{Error messages should say what to do}

The startup test does not say "assertion failed". It names the modules that were
loaded and points at the file that shows the correct pattern. Writing that
message took an extra minute and will save an hour for whoever hits it.

\subsection{Small representation changes remove whole bug families}

Month ranges kept producing boundary errors until a year and a month became one
integer. After that, a range is two comparisons and there are no special cases
left to get wrong. When conditions keep multiplying, the representation is
usually the problem.

\section{What would be done differently}

\begin{itemize}
  \item \textbf{The content security policy still allows inline scripts},
        because chart initialisation lives in the templates. Moving that code to
        static files would let the policy drop \code{'unsafe-inline'}.
  \item \textbf{The Trade Republic parser depends on a PDF layout} and will
        break when the format changes. A statement fixture per layout version
        would turn that break into a failing test instead of a silent
        misreading.
  \item \textbf{The login rate limiter lives in memory}, so a restart clears it.
        Correct for one process, wrong the moment a second one is added.
  \item \textbf{Migrations were reconstructed rather than translated} for the
        public repository. That was right here, and it would not be an option
        for a project with real deployments to upgrade.
\end{itemize}

\section{Closing}

The project began as five desktop scripts that only worked with one computer
switched on. It ended as a single application that runs on a small server,
answers from a phone, keeps its own history, reports on itself, and survives a
free hosting plan without supervision.

The technical work was not the hard part. The hard parts were deciding what not
to build, being honest about what fails and how, and making the thing keep
working when nobody is watching it. Those are the parts this report tried to
make explicit.
